%! TEX root = master.tex
\chapter{2}{Power and Taylor Series}

Every function analytic on a disc is the sum of its Taylor series there, and
power series can be differentiated, multiplied and compared term by term. Each
section states one of these facts and then puts it to work.

\section{Radius of convergence}\label{sec:radius}

%! TEX root = ../master.tex
\begin{theorem}[Cauchy--Hadamard]\label{thm:cauchy-hadamard}
  Let $\sum a_n (z - z_0)^n$ be a power series and define $R \in [0, \infty]$ by
  \[
    \frac{1}{R} = \limsup_{n \to \infty} \abs{a_n}^{1/n},
  \]
  with $1/0 = \infty$ and $1/\infty = 0$. The series converges absolutely for
  $\abs{z - z_0} < R$ and diverges for $\abs{z - z_0} > R$. We call $R$ the
  \emph{radius of convergence}.
  \thmsource{\cite[ch.~5]{saff-snider}}
\end{theorem}

%! TEX root = ../master.tex
\begin{theorem}[Ratio test]\label{thm:ratio-test}
  If $a_n \neq 0$ for all large $n$ and
  $L = \lim_{n \to \infty} \abs{a_{n+1}/a_n}$ exists in $[0, \infty]$, then the
  radius of convergence of $\sum a_n (z - z_0)^n$ is $R = 1/L$.
  \thmsource{\cite[ch.~5]{saff-snider}}
\end{theorem}


\section{Termwise differentiation}\label{sec:termwise}

%! TEX root = ../master.tex
\begin{theorem}[Termwise differentiation]\label{thm:termwise}
  Suppose $f(z) = \sum_{n=0}^{\infty} a_n (z - z_0)^n$ has radius of
  convergence $R > 0$. Then $f$ is analytic on $\abs{z - z_0} < R$, and
  \[
    f'(z) = \sum_{n=1}^{\infty} n a_n (z - z_0)^{n-1},
  \]
  where the differentiated series also has radius of convergence $R$.
  \thmsource{\cite[ch.~5]{saff-snider}}
\end{theorem}

\begin{proof}
  Translating, we may take $z_0 = 0$.

  \emph{Same radius.} Since $n^{1/n} \to 1$,
  \[
    \limsup_{n \to \infty} \abs{n a_n}^{1/n}
    = \limsup_{n \to \infty} \abs{a_n}^{1/n},
  \]
  so $\sum n a_n z^{n-1}$ has radius $R$ by \cref{thm:cauchy-hadamard}.
  Similarly, $\sum n(n-1) \abs{a_n} r^{n-2}$ converges for every $0 \leq r < R$.

  \emph{The derivative.} Fix $\abs{z} < R$, choose $r$ with $\abs{z} < r < R$,
  and let $g(z) = \sum n a_n z^{n-1}$. For $\abs{w} < r$ with $w \neq z$,
  \[
    \frac{f(w) - f(z)}{w - z} - g(z)
    = \sum_{n=2}^{\infty} a_n \left( \frac{w^n - z^n}{w - z} - n z^{n-1} \right).
  \]
  Expanding $\frac{w^n - z^n}{w - z} = \sum_{j=0}^{n-1} w^j z^{n-1-j}$, the
  bracket equals $\sum_{j=1}^{n-1} z^{n-1-j} (w^j - z^j)$. Since
  $\abs{w^j - z^j} \leq j r^{j-1} \abs{w - z}$, the bracket is at most
  \[
    \abs{w - z} \, r^{n-2} \sum_{j=1}^{n-1} j
    = \abs{w - z} \, \frac{n(n-1)}{2} \, r^{n-2}
  \]
  in modulus. Hence
  \[
    \abs*{\frac{f(w) - f(z)}{w - z} - g(z)}
    \leq \frac{\abs{w - z}}{2} \sum_{n=2}^{\infty} n(n-1) \abs{a_n} r^{n-2}
    \longrightarrow 0
    \quad \text{as } w \to z.
  \]
  So $f'(z) = g(z)$ exists at every point of the disc. That is, $f$ is
  analytic there.
\end{proof}


\begin{exercise}\label{ex:sum-n-cubed}
  Find the radius of convergence and a closed form for
  $f(z) = \sum_{n=1}^{\infty} n^3 z^n$. Use it to evaluate
  $\sum_{n=1}^{\infty} n^3 / 3^n$ and $\sum_{n=3}^{\infty} n^3 / 2^n$.
\end{exercise}

\begin{solution}
  \emph{Radius.} With $a_n = n^3$,
  \[
    \abs*{\frac{a_{n+1}}{a_n}} = \left(1 + \frac{1}{n}\right)^3 \to 1,
  \]
  so $R = 1$ by \cref{thm:ratio-test}.

  \emph{Closed form.} Start from the geometric series
  \[
    \frac{1}{1 - z} = \sum_{n=0}^{\infty} z^n, \qquad \abs{z} < 1.
  \]
  Differentiating term by term (\cref{thm:termwise}) and multiplying by $z$
  replaces each coefficient $a_n$ by $n a_n$ without changing the disc.
  Doing this three times gives
  \begin{align*}
    \frac{z}{(1-z)^2}             &= \sum_{n=1}^{\infty} n z^n, \\
    \frac{z(1+z)}{(1-z)^3}        &= \sum_{n=1}^{\infty} n^2 z^n, \\
    \frac{z(1 + 4z + z^2)}{(1-z)^4} &= \sum_{n=1}^{\infty} n^3 z^n.
  \end{align*}
  For instance, the last line comes from the quotient rule:
  \[
    \frac{\dd}{\dd z} \frac{z + z^2}{(1-z)^3}
    = \frac{(1 + 2z)(1 - z) + 3(z + z^2)}{(1-z)^4}
    = \frac{1 + 4z + z^2}{(1-z)^4}.
  \]

  \emph{The sums.} Both points lie inside $\abs{z} < 1$. At $z = 1/3$,
  \[
    \sum_{n=1}^{\infty} \frac{n^3}{3^n}
    = \frac{\frac13 \cdot \frac{22}{9}}{\left(\frac23\right)^4}
    = \frac{22}{27} \cdot \frac{81}{16}
    = \frac{33}{8}.
  \]
  At $z = 1/2$ the full sum is
  $\frac{\frac12 \cdot \frac{13}{4}}{\left(\frac12\right)^4} = 26$, and removing
  the terms $n = 1, 2$, which are $\frac12$ and $2$, leaves
  \[
    \sum_{n=3}^{\infty} \frac{n^3}{2^n} = 26 - \frac52 = \frac{47}{2}. \qedhere
  \]
\end{solution}

\begin{remark}
  The operator $\vartheta = z \frac{\dd}{\dd z}$ sends $z^n$ to $n z^n$, so
  $f = \vartheta^3 \bigl(\frac{1}{1-z}\bigr)$. The same trick sums
  $\sum p(n) z^n$ for any polynomial $p$.
\end{remark}

\begin{exercise}\label{ex:bessel}
  Let
  \[
    f(z) = \sum_{n=0}^{\infty} \frac{(-1)^n z^{2n}}{4^n (n!)^2}.
  \]
  Show that $f$ is entire and that
  $f''(z) + \frac{1}{z} f'(z) + f(z) = 0$ for $z \neq 0$.
\end{exercise}

\begin{proof}
  \emph{$f$ is entire.} Fix $z \neq 0$ and let $c_n$ be the $n$th term. Then
  \[
    \abs*{\frac{c_{n+1}}{c_n}} = \frac{\abs{z}^2}{4(n+1)^2} \to 0,
  \]
  so the series converges absolutely for every $z$. Its radius of convergence
  is $\infty$, and $f$ is entire by \cref{thm:termwise}.

  \emph{Derivatives.} By \cref{thm:termwise}, twice,
  \[
    f'(z) = \sum_{n=1}^{\infty} \frac{(-1)^n \, 2n \, z^{2n-1}}{4^n (n!)^2},
    \qquad
    f''(z) = \sum_{n=1}^{\infty}
      \frac{(-1)^n \, 2n(2n-1) \, z^{2n-2}}{4^n (n!)^2}.
  \]
  Every term of $f'$ contains a factor $z$, so
  \[
    \frac{1}{z} f'(z)
    = \sum_{n=1}^{\infty} \frac{(-1)^n \, 2n \, z^{2n-2}}{4^n (n!)^2}
  \]
  is again a convergent power series. In fact it is entire, so the
  singularity of $f'(z)/z$ at the origin is only apparent.

  \emph{The equation.} Adding the last two series term by term,
  \[
    f''(z) + \frac{1}{z} f'(z)
    = \sum_{n=1}^{\infty} \frac{(-1)^n \, 4n^2 \, z^{2n-2}}{4^n (n!)^2}
    = \sum_{n=1}^{\infty} \frac{(-1)^n z^{2n-2}}{4^{n-1} ((n-1)!)^2},
  \]
  since $4^n = 4 \cdot 4^{n-1}$ and $(n!)^2 = n^2 ((n-1)!)^2$. Put $k = n - 1$:
  \[
    f''(z) + \frac{1}{z} f'(z)
    = \sum_{k=0}^{\infty} \frac{(-1)^{k+1} z^{2k}}{4^k (k!)^2}
    = -f(z). \qedhere
  \]
\end{proof}

\begin{remark}
  This $f$ is the Bessel function $J_0$. Multiplying through by $z$ gives
  Bessel's equation of order zero, $z f'' + f' + z f = 0$, which then holds
  at $z = 0$ as well.
\end{remark}

\section{Uniqueness of coefficients}\label{sec:uniqueness}

%! TEX root = ../master.tex
\begin{theorem}[Uniqueness of power series]\label{thm:uniqueness}
  If $\sum a_n (z - z_0)^n$ converges to $f(z)$ on some disc
  $\abs{z - z_0} < r$, then
  \[
    a_n = \frac{f^{(n)}(z_0)}{n!} \qquad (n \geq 0).
  \]
  In particular, two power series with the same sum near $z_0$ have the same
  coefficients.
  \thmsource{\cite[ch.~5]{saff-snider}}
\end{theorem}

\begin{proof}
  Applying \cref{thm:termwise} $k$ times,
  \[
    f^{(k)}(z) = \sum_{n=k}^{\infty} \frac{n!}{(n-k)!} \, a_n (z - z_0)^{n-k},
    \qquad \abs{z - z_0} < r.
  \]
  Putting $z = z_0$ kills every term except $n = k$, so $f^{(k)}(z_0) = k! \, a_k$.
  If two series have the same sum $f$ near $z_0$, both sets of coefficients
  equal $f^{(n)}(z_0)/n!$.
\end{proof}


\begin{exercise}\label{ex:lacunary}
  Let $f(z) = \sum_{n=0}^{\infty} z^{2^n} = z + z^2 + z^4 + z^8 + \dotsb$.
  Show that $f(z) = z + f(z^2)$, and find every power series $g$ that
  satisfies $g(z) = z + g(z^2)$ near the origin.
\end{exercise}

\begin{solution}
  \emph{$f$ satisfies the equation.} The coefficients are $0$ or $1$, with
  infinitely many $1$s, so the radius of convergence is $1$. For $\abs{z} < 1$,
  \[
    z + f(z^2) = z + \sum_{n=0}^{\infty} z^{2^{n+1}}
      = z^{2^0} + \sum_{n=1}^{\infty} z^{2^n} = f(z).
  \]

  \emph{All solutions.} Suppose $g(z) = \sum a_n z^n$ converges near $0$ and
  $g(z) = z + g(z^2)$. Then
  \[
    a_0 + a_1 z + a_2 z^2 + a_3 z^3 + a_4 z^4 + \dotsb
    = a_0 + z + a_1 z^2 + 0 z^3 + a_2 z^4 + \dotsb,
  \]
  and by \cref{thm:uniqueness}
  \[
    a_0 \text{ is free}, \qquad a_1 = 1, \qquad
    a_{2k} = a_k \ (k \geq 1), \qquad a_{2k+1} = 0 \ (k \geq 1).
  \]
  Write $n \geq 1$ as $n = 2^m q$ with $q$ odd. Halving $m$ times gives
  $a_n = a_q$, which is $1$ if $q = 1$ and $0$ otherwise. So $a_n = 1$
  exactly when $n$ is a power of $2$, and
  \[
    g(z) = c + f(z), \qquad c \in \C.
  \]
  Conversely, each such $g$ works, because the constant appears on both sides
  of $g(z) = z + g(z^2)$. The solutions are the translates of $f$ by
  constants. \qedhere
\end{solution}

\begin{remark}
  The gaps in the exponents $2^n$ make $f$ impossible to continue past the
  unit circle: every point of $\abs{z} = 1$ is a singularity. This is a
  \emph{natural boundary}, proved as a separate graduate proof.
  % TODO(review): \cref the natural-boundary proof once it exists.
\end{remark}

\section{Products and Taylor's theorem}\label{sec:products}

%! TEX root = ../master.tex
\begin{theorem}[Cauchy product]\label{thm:cauchy-product}
  If $f(z) = \sum a_n (z - z_0)^n$ and $g(z) = \sum b_n (z - z_0)^n$ both
  converge on $\abs{z - z_0} < r$, then on the same disc
  \[
    f(z) g(z) = \sum_{n=0}^{\infty} c_n (z - z_0)^n,
    \qquad c_n = \sum_{k=0}^{n} a_k b_{n-k}.
  \]
  \thmsource{\cite[ch.~5]{saff-snider}}
\end{theorem}

%! TEX root = ../master.tex
\begin{theorem}[Taylor's theorem]\label{thm:taylor}
  If $f$ is analytic on $\abs{z - z_0} < \rho$, then
  \[
    f(z) = \sum_{n=0}^{\infty} \frac{f^{(n)}(z_0)}{n!} (z - z_0)^n
  \]
  for every $z$ in that disc. In particular, the radius of convergence of the
  Taylor series is at least the distance from $z_0$ to the nearest point where
  $f$ fails to be analytic.
  \thmsource{\cite[ch.~5]{saff-snider}}
\end{theorem}


\begin{exercise}\label{ex:secant}
  Find the first four nonzero terms of the Maclaurin series of $1/\cos z$.
  Justify each step.
\end{exercise}

\begin{solution}
  \emph{A series exists.} $\cos z$ is entire and $\cos 0 = 1 \neq 0$, so
  $1/\cos z$ is analytic near $0$. By Taylor's theorem it has a Maclaurin
  series $\sum b_n z^n$.

  \emph{Multiply out.} Write
  $\cos z = 1 - \frac{z^2}{2} + \frac{z^4}{24} - \frac{z^6}{720} + \dotsb$.
  Since $\cos z \cdot \frac{1}{\cos z} = 1$, \cref{thm:cauchy-product} gives
  \[
    \Bigl(1 - \tfrac12 z^2 + \tfrac1{24} z^4 - \tfrac1{720} z^6 + \dotsb\Bigr)
    \bigl(b_0 + b_1 z + b_2 z^2 + \dotsb\bigr)
    = 1 + 0z + 0z^2 + \dotsb.
  \]

  \emph{Compare coefficients.} By \cref{thm:uniqueness} the coefficients of
  each power of $z$ agree on both sides:
  \begin{alignat*}{3}
    & z^0\colon\quad & b_0 &= 1, \\
    & z^2\colon & b_2 - \tfrac12 b_0 &= 0, &\qquad& b_2 = \tfrac12, \\
    & z^4\colon & b_4 - \tfrac12 b_2 + \tfrac1{24} b_0 &= 0, && b_4 = \tfrac{5}{24}, \\
    & z^6\colon & b_6 - \tfrac12 b_4 + \tfrac1{24} b_2 - \tfrac1{720} b_0 &= 0, &&
      b_6 = \tfrac{61}{720}.
  \end{alignat*}
  The odd equations read $b_1 = 0$, $b_3 = \frac12 b_1 = 0$, and so on, as
  they must, since $1/\cos z$ is even. Hence
  \[
    \frac{1}{\cos z} = 1 + \frac{z^2}{2} + \frac{5 z^4}{24}
      + \frac{61 z^6}{720} + \dotsb. \qedhere
  \]
\end{solution}

\begin{remark}
  The zeros of $\cos z$ nearest the origin are $\pm \pi/2$, so this series
  converges exactly on $\abs{z} < \pi/2$. The numerators $1, 1, 5, 61$ are the
  Euler (secant) numbers.
\end{remark}

\begin{exercise}\label{ex:z-fourth-about-one}
  Expand $p(z) = z^4$ in powers of $z - 1$.
\end{exercise}

\begin{solution}
  $p$ is entire, so by \cref{thm:taylor} its Taylor series about $1$ converges
  to it everywhere. The derivatives at $1$ are
  $p(1) = 1$, $p'(1) = 4$, $p''(1) = 12$, $p'''(1) = 24$, $p^{(4)}(1) = 24$, and
  $p^{(j)} \equiv 0$ for $j \geq 5$. Dividing by $j!$,
  \[
    z^4 = 1 + 4(z - 1) + 6(z - 1)^2 + 4(z - 1)^3 + (z - 1)^4. \qedhere
  \]
\end{solution}

\begin{remark}
  The coefficients are the binomial coefficients $\binom{4}{j}$, as they must
  be: $z^4 = \bigl(1 + (z - 1)\bigr)^4$.
\end{remark}
